% !TEX root = ../print.tex

\chapter{Gradings}

\section{Grading types}

The morphisms of an \Ainf{}-category are graded by an abelian group $\beta$ of
degrees. The \Ainf{}-axioms use exactly two further pieces of structure on
$\beta$: integer degree shifts, so that the $n$-ary operation $m_n$ can have
degree $2-n$, and a Koszul sign attached to every degree. No degree is ever
multiplied, so $\beta$ need not be a ring. Besides the standard gradings
$\mathbb{Z}$ and $\mathbb{Z}/2$, this allows bigradings such as
$\mathbb{Z}\times\mathbb{Z}$ or $\mathbb{Z}\times\mathbb{Z}/2$, in which $m_n$
has degree $(2-n,0)$.

\begin{definition}[Grading type]
\lean{AInfinityTheory.GradingType}
\label{AInfinityTheory.GradingType}
\leanok
A grading type is an abelian group $\beta$ together with group homomorphisms
\[
  \operatorname{shift} : \mathbb{Z} \to \beta,
  \qquad
  \operatorname{sign} : \beta \to \{\pm 1\},
\]
such that $\operatorname{sign}(\operatorname{shift}(1)) = -1$; that is, the
differential, which has degree $\operatorname{shift}(1)$, is odd. For a degree
$d \in \beta$ we write $(-1)^d := \operatorname{sign}(d)$ for its Koszul sign.
\end{definition}

\begin{lemma}[Sign of an integer shift]
\lean{AInfinityTheory.sign_shift}
\label{AInfinityTheory.sign_shift}
\leanok
\uses{AInfinityTheory.GradingType}
For every integer $n$ one has $(-1)^{\operatorname{shift}(n)} = (-1)^n$.
\end{lemma}
\begin{proof}
\leanok
Since $\operatorname{shift}(n) = n \cdot \operatorname{shift}(1)$ and
$\operatorname{sign}$ is a homomorphism,
$(-1)^{\operatorname{shift}(n)} = \operatorname{sign}(\operatorname{shift}(1))^n = (-1)^n$.
\end{proof}

\begin{definition}[Standard gradings]
\lean{AInfinityTheory.GradingType.int, AInfinityTheory.GradingType.zmodTwo}
\label{AInfinityTheory.GradingType.standard}
\leanok
\uses{AInfinityTheory.GradingType}
The integers $\mathbb{Z}$ form a grading type with $\operatorname{shift}$ the
identity and $\operatorname{sign}(n) = (-1)^n$. The parities $\mathbb{Z}/2$
form a grading type with $\operatorname{shift}$ the reduction modulo $2$ and
$\operatorname{sign}(p) = (-1)^p$.
\end{definition}

\begin{definition}[Product gradings]
\lean{AInfinityTheory.GradingType.prodFst, AInfinityTheory.GradingType.prodTotal}
\label{AInfinityTheory.GradingType.prod}
\leanok
\uses{AInfinityTheory.GradingType}
Let $\beta$ be a grading type and $\gamma$ an abelian group. The product
$\beta \times \gamma$ is a grading type with
$\operatorname{shift}(n) = (\operatorname{shift}(n), 0)$ in two ways: either
the sign is read off the first factor alone, $\operatorname{sign}(b,c) = (-1)^b$,
or, when $\gamma$ is a grading type as well, the sign is the total sign
$\operatorname{sign}(b,c) = (-1)^b\,(-1)^c$. Neither choice is canonical, so
both are constructions to be activated locally rather than instances.
\end{definition}

\section{Graded modules and quivers}

The graded-module and graded-quiver definitions below do not need a grading
type; the type of degrees is arbitrary there. The structure of a
grading type is only imposed by the degree arithmetic and Koszul signs of the
\Ainf{}-axioms.

\begin{definition}[Graded $R$-module]
\lean{AInfinityTheory.GradedRModule}
\label{AInfinityTheory.GradedRModule}
\leanok
For a commutative ring $R$, a graded $R$-module is a $\beta$-indexed family of
$R$-modules, implemented as a graded object in \texttt{ModuleCat R}.
\end{definition}

\begin{definition}[$R$-linear graded quiver]
\lean{AInfinityTheory.RLinearGradedQuiver}
\label{AInfinityTheory.RLinearGradedQuiver}
\leanok
\uses{AInfinityTheory.GradedRModule}
An $R$-linear graded quiver on an object type $\mathrm{Obj}$ assigns to each pair
of objects $(X,Y)$ a graded $R$-module of morphisms from $X$ to $Y$.
\end{definition}
